\scanpage{13}
\mspage{9}
% source: top, continuation
Mais c'est à l'ens des supports \uline{pour $\unsure{M|\mathscr M}$} qu'on s'intéresse en fait!

% source: upper part, figures and polygons
Nous allons à présent travailler avec des figures et polygones finis de $\widehat M$ (ou $\widehat M$-polygones, ou polygones \unsure{généralisés}), i.e. les parties finies tot. \unsure{ordonnées} (ou drapeaux) de $\widehat M,\leq$. Pour une telle partie $\unsure{\Phi|\varphi}$, on introduit
\[
T^{\Phi}=\delta\Phi\in\operatorname{Drap}^{\unsure{\times|*}}(\unsure{L|\mathscr L}),
\]
qui est vide ssi $\Phi$ est vide.
% source: left margin beside definition of T
% slanted
\margin{\unsure{On va par la} suite se borner aux non vides.}

% source: middle, interstitial strata
Ainsi on introduit les \uline{strates interstitielles} (\unsure{généralisées|généralisation}) d'un drapeau
\[
T=\{t_{\unsure{1|l}}<t_2<\cdots<t_n\}
\]
de $L$, ainsi: \deleted{\unsure{Si on prend}}
\[
[I_{<t_1}],\ t_1,\ I_{t_1,t_2},\ t_2,\ I_{t_2,t_3},\ \cdots,\ t_n,[I_{>t_n}]
\]
% source: lower-middle, qualifying sentence and insertion
où le terme $I_{<t_1}$ ne figure que si $t_1$ non minimal\inserted{(sinon, il n'existe \unsure{même} pas)}, et le terme $I_{>t_n}$ ne figure que si $t_n$ est non maximal (\unsure{idem}\ldots).

% source: lower, largest polygon
Ce sont les éléments de le plus grand $\widehat M$-polygone possible $\widehat\Phi$, telle que $\delta\widehat\Phi=T$.

% source: bottom, data defining Phi
La donnée de \inserted{\unsure{le} $\widehat M$-polygone} $\Phi\ne\varnothing$ revient à celle du drapeau $T$ plus un ensemble fini\inserted{$\Phi$,} des strates interstitielles (\unsure{généralisées}) ayant la propriété suivante:
\[
\forall t\in T,\quad\text{on a}\quad t\in\Phi,\quad\uline{\text{ou}}\quad t\triangleleft X\quad\text{avec }X\in\Phi.
\]
% archive sheet number: 12
